\Article{10}
Drittens soll die Oberfläche einer Kugel mit dem Radius \(=a\) in der Ebene abgebildet werden. Wir setzen hier
\[
\begin{aligned}
x&=a\cos t\mathord{.}\sin u\\
y&=a\sin t\mathord{.}\sin u\\
z&=a\cos u
\end{aligned}
\]
\noindent woraus wir erhalten
\[
\omega=aa\sin u^2\,\mathrm{d}t^2+aa\,\mathrm{d}u^2
\]
\noindent Aus der Differentialformel \(\omega=0\) folgt daher
\[
\mathrm{d}t\mp i\mathord{.}\frac{\mathrm{d}u}{\sin u}=0
\]
\noindent und durch Integration ergibt sich
\[
t\pm i\log\mathrm{cotang}\tfrac{1}{2}u=\text{Const.}
\]
\par\indent\SourceLeaf{209}%
Wenn wir also mit dem Zeichen \(f\) wieder eine beliebige Funktion bezeichnen, müssen wir \(X\) dem Realteil und \(iY\) dem Imaginärteil von
\[
f(t+i\log\mathrm{cotang}\tfrac{1}{2}u)
\]
\noindent gleichsetzen. Wir wollen einige Sonderfälle dieser allgemeinen Lösung anführen.

Wählt man für \(f\) eine lineare Funktion, indem man \(f\upsilon=k\upsilon\) setzt, so erhält man
\[
X=kt,\quad Y=k\log\mathrm{cotang}\tfrac{1}{2}u
\]
\noindent Auf die Erde angewandt, stimmt dies offensichtlich mit der \textsc{Mercator}-Projektion überein, wenn \(t\) die geografische Länge und \(90^0-u\) die Breite bezeichnet. Für das Vergrößerungsverhältnis ergeben die Formeln aus Artikel 7 hier
\[
m=\frac{k}{a\sin u}
\]

Wählt man für \(f\) eine imaginäre Exponentialfunktion und zunächst die einfachste \(f\upsilon=ke^{i\upsilon}\), so erhält man
\[
f(t+i\log\mathrm{cotang}\tfrac{1}{2}u)=ke^{\log\mathrm{tang}\frac{1}{2}u+it}=k\,\mathrm{tang}\tfrac{1}{2}u(\cos t+i\sin t)
\]
\noindent und
\[
X=k\,\mathrm{tang}\tfrac{1}{2}u\mathord{.}\cos t,\quad
Y=k\,\mathrm{tang}\tfrac{1}{2}u\mathord{.}\sin t
\]
\noindent Das ist, wie man leicht sieht, die stereografische Polarprojektion.

Setzt man allgemeiner \(f\upsilon=ke^{i\lambda\upsilon}\), so erhält man
\[
X=k\,\mathrm{tang}\tfrac{1}{2}u^{\lambda}\mathord{.}\cos\lambda t,\quad
Y=k\,\mathrm{tang}\tfrac{1}{2}u^{\lambda}\mathord{.}\sin\lambda t
\]
\noindent Für das Vergrößerungsverhältnis erhalten wir hier
\[
n=aa\sin u^2,\quad N=1,\quad\varphi\upsilon=i\lambda ke^{i\lambda\upsilon}
\]
\noindent und daraus
\[
m=\frac{\lambda k\,\mathrm{tang}\frac{1}{2}u^{\lambda}}{a\sin u}
\]

Man sieht, dass die Abbildung aller Punkte mit konstantem \(u\) auf einem einzigen Kreis liegt und die Abbildung aller Punkte mit konstantem \(t\) auf einer einzigen Geraden. Die Kreise zu den verschiedenen Werten von \(u\) sind konzentrisch. Dies ergibt eine sehr zweckmäßige Kartenprojektion, wenn nur ein Teil der Kugeloberfläche abzubilden ist. Dann wählt man \(\lambda\) am besten so, dass das Vergrößerungsverhältnis für die
\SourceLeaf{210}
äußersten Werte von \(u\) gleich groß wird und dadurch zur Mitte hin seinen kleinsten Wert erreicht. Sind diese äußersten Werte von \(u\) durch \(u^0\) und \(u'\) gegeben, so muss man folglich setzen
\[
\lambda=\frac{\log\sin u'-\log\sin u^0}{\log\mathrm{tang}\frac{1}{2}u'-\log\mathrm{tang}\frac{1}{2}u^0}
\]
\noindent Die Blätter Nr. 19---26 der Sternkarten von Professor \textsc{Harding} sind nach dieser Projektion gezeichnet.

\Article{11}
Die allgemeine Lösung für das im vorigen Artikel behandelte Beispiel lässt sich noch in einer anderen Form angeben. Wegen ihrer Eleganz glauben wir, sie hier hinzufügen zu sollen.

Aus dem in Artikel 6 Dargelegten folgt, da
\[
\mathrm{tang}\tfrac{1}{2}u(\cos t+i\sin t)
\]
\noindent eine Funktion von
\[
t+i\log\mathrm{cotang}\tfrac{1}{2}u
\]
\noindent ist und
\[
\mathrm{tang}\tfrac{1}{2}u(\cos t+i\sin t)=\frac{\sin u\cos t+i\sin u\sin t}{1+\cos u}=\frac{x+iy}{a+z}
\]
\noindent dass sich die allgemeine Lösung auch durch
\[
X+iY=f\frac{x+iy}{a+z},\quad X-iY=f'\frac{x-iy}{a+z}
\]
\noindent darstellen lässt. Das heißt, \(X\) muss dem Realteil und \(iY\) dem Imaginärteil von \(f\frac{x+iy}{a+z}\) gleichgesetzt werden, wobei \(f\) eine beliebige Funktion bezeichnet. Anstelle von \(f\frac{x+iy}{a+z}\) kann man, wie man leicht sieht, auch eine beliebige Funktion von \(\frac{y+iz}{a+x}\) oder von \(\frac{z+ix}{a+y}\) nehmen.

